% !TeX spellcheck = pt_BR
% ============================================================
% Tabela de constantes fisicas
% Valores adotados ao longo do livro
% ============================================================

\chapter*{Tabela de Constantes}
\addcontentsline{toc}{chapter}{Tabela de Constantes}

Ao longo do livro, salvo indicação contrária, adotam-se os seguintes
valores para as constantes físicas:

\vspace{1em}

\begin{center}
\renewcommand{\arraystretch}{1.4}
\begin{tabular}{l c >{\itshape}l}
\toprule
\textbf{Grandeza} & \textbf{Valor adotado} & \textnormal{\textbf{Unidade SI}} \\
\midrule
Aceleração da gravidade (Terra)              & $g = 9{,}8$    & m/s$^2$ \\
Aceleração da gravidade (valor didático)     & $g = 10$       & m/s$^2$ \\
Constante gravitacional universal            & $G = 6{,}67 \times 10^{-11}$ & N$\cdot$m$^2$/kg$^2$ \\
Velocidade da luz no vácuo                   & $c = 3 \times 10^{8}$        & m/s \\
Densidade da água (a 4$^\circ$C)             & $\rho = 1{,}0 \times 10^{3}$ & kg/m$^3$ \\
Densidade do ar (nível do mar)               & $\rho_\text{ar} = 1{,}3$     & kg/m$^3$ \\
Pressão atmosférica (nível do mar)           & $P_0 = 1{,}01 \times 10^{5}$ & Pa \\
\midrule
\multicolumn{3}{l}{\itshape Dados da Terra} \\
Massa da Terra                               & $M_T = 6{,}0 \times 10^{24}$ & kg \\
Raio médio da Terra                          & $R_T = 6{,}4 \times 10^{6}$  & m \\
Distância Terra--Sol (média)                 & $d_{TS} = 1{,}5 \times 10^{11}$ & m \\
\midrule
\multicolumn{3}{l}{\itshape Convenções de arredondamento} \\
\multicolumn{3}{l}{\small Quando o problema é do tipo ENEM/vestibular, $g = 10$ m/s$^2$.} \\
\multicolumn{3}{l}{\small Em demonstrações teóricas e olímpiadas, $g = 9{,}8$ m/s$^2$.} \\
\bottomrule
\end{tabular}
\end{center}

\vspace{1em}

\noindent\textit{Nota: os valores acima são aproximações usadas para fins
didáticos. Para precisão científica, consulte a base de dados do
CODATA mais recente.}

\clearpage
